\section{Related work}
\label{sec:related}

\paragraph{Sequence-level distributional control}
Posing controlled generation as a moment constraint and solving for the KL-minimal
distribution yields an exponential-family optimum, and distributional policy gradient
trains a model toward it by importance-weighted sampling against that unnormalized
target~\citep{khalifa2021distributional}. The formulation is attractive because the optimum
is characterized in closed form up to a normalizer and the multiplier has a direct reading
as the price of the constraint. Its practical difficulty is that the target is defined over
complete sequences, so the training signal arrives only once a sequence terminates and the
estimator's variance grows with the tilt strength. Where the resulting distribution is
evaluated, it is evaluated on whether the constrained expectation is met, which is the
quantity the formulation optimizes and not the quantity a user inspecting samples
experiences.

\paragraph{Expectation versus per-instance constraints}
The distinction our concentration result turns on is the one constrained Markov
decision processes make between constraining an expectation and constraining every
trajectory~\citep{altman1999cmdp}, refined by state-wise
formulations~\citep{zhao2023scpo} and by chance-constrained
programming~\citep{charnes1959chance}. The precedent closest in mechanism is the safety
layer, which projects a proposed action onto a feasible set at every
step~\citep{pham2018optlayer}. We adopt that vocabulary: the sequence-level tilt is an
expectation constraint, and what our projection delivers is closer to a per-step
feasibility guarantee whose aggregate happens to meet the expectation.

\paragraph{Per-step constrained decoding}
Lexically constrained decoding enforces predicate-logic constraints during search
without retraining~\citep{lu2021neurologic}, with lookahead heuristics added to reduce
myopia~\citep{lu2022neurologic}. GeLaTo instead distills a tractable generative model to
supply exact conditional guidance~\citep{zhang2023gelato}, and sequential Monte Carlo
methods reweight partial sequences against a target
potential~\citep{zhao2024twisted, loula2025syntactic}. All three act at the per-step locus but treat the
constraint as a filter over candidates or a reweighting of them, so none yields a
distribution that is itself the solution of a constrained problem, and none is
differentiable with respect to the constraint. Ahmed
\etal~\citep{ahmed2023pseudosemantic} are the nearest training-time neighbour, relaxing
a logical constraint into a pseudo-semantic loss over a local approximation, the
same shape of move as reducing a sequence-level expectation to a per-step one, though
approximate in the constraint rather than exact in the support.

\paragraph{Differentiable optimization layers}
Embedding an argmin in a network and backpropagating through its solution is established
for quadratic programs, where differentiating the Karush--Kuhn--Tucker conditions gives the
gradient in closed form~\citep{amos2017optnet}; for general convex programs, where the same
idea is automated by rewriting the problem in a canonical cone
form~\citep{agrawal2019cvxpylayers}; and for hard constraints, where a completion-plus-
correction scheme restores feasibility that a learned map would otherwise
violate~\citep{donti2021dc3}. What these share is a cost profile: a generic solver is
called in the forward pass and its implicit derivative in the backward pass, so the layer
is affordable per example but not per token. None of them has been instantiated as a
divergence projection over a next-token distribution, where the problem is small and
structured enough that a generic solver is the wrong tool.

\paragraph{Crystal structure generation}
Generative models for inorganic crystals include diffusion over invariant
representations~\citep{xie2022cdvae, jiao2023diffcsp}, autoregressive language
models over text serializations~\citep{antunes2024crystallm, gruver2024finetuned}, and
large-scale diffusion trained for property targeting~\citep{zeni2025mattergen}. We use
the same MP-20 subset of the Materials Project~\citep{jain2013materialsproject} and a
text serialization in the spirit of the autoregressive line, because a serialization is
what makes a per-step budget definable at all. That the class is bounded this way is a real
constraint on the method rather than an implementation detail.

%%%%%%%%%%%%%%%%%%%%%%%%%%%%%%%%%%%%%%%%%%%%%%%%%%%%%%%%%%%%%%%%%%%%%%%%%%%%%%%%
